\documentclass[11pt]{amsart}
\usepackage[localtheorems]{raeez-math-template}

\providecommand{\C}{\mathbb{C}}
\providecommand{\R}{\mathbb{R}}
\providecommand{\Z}{\mathbb{Z}}
\providecommand{\N}{\mathbb{N}}
\providecommand{\Q}{\mathbb{Q}}
\providecommand{\Pp}{\mathbb{P}}
\providecommand{\hb}{\hbar}
\providecommand{\eps}{\varepsilon}
\providecommand{\Sym}{\mathrm{Sym}}
\providecommand{\Hilb}{\mathrm{Hilb}}
\providecommand{\Coh}{\mathrm{Coh}}
\providecommand{\DT}{\mathrm{DT}}
\providecommand{\PT}{\mathrm{PT}}
\providecommand{\Aut}{\mathrm{Aut}}
\providecommand{\Tot}{\mathrm{Tot}}
\providecommand{\Obs}{\mathrm{Obs}}
\providecommand{\Perf}{\mathrm{Perf}}
\providecommand{\CoHA}{\mathrm{CoHA}}
\providecommand{\Rep}{\mathrm{Rep}}
\providecommand{\CR}{\mathrm{CR}}
\providecommand{\Gr}{\mathrm{Gr}}
\providecommand{\cO}{\mathcal{O}}
\providecommand{\cM}{\mathcal{M}}
\providecommand{\cP}{\mathcal{P}}
\providecommand{\cF}{\mathcal{F}}
\providecommand{\cE}{\mathcal{E}}
\providecommand{\cH}{\mathcal{H}}
\providecommand{\cZ}{\mathcal{Z}}
\providecommand{\kch}{\kappa_{\mathrm{ch}}}
\providecommand{\kBKM}{\kappa_{\mathrm{BKM}}}
\providecommand{\kcat}{\kappa_{\mathrm{cat}}}
\providecommand{\kfib}{\kappa_{\mathrm{fiber}}}
\providecommand{\Meff}{\cM^{+}_{\mathrm{eff}}}
\providecommand{\Yp}{Y^{+}}
\providecommand{\Yph}{Y^{+}(\widehat{\mathfrak{gl}}_1)}
\providecommand{\gDelta}{\mathfrak{g}_{\Delta_5}}
\providecommand{\Phion}{\phi_{0,1}}
\providecommand{\SpCh}{\mathrm{Sp}^{\mathrm{ch}}}
\providecommand{\PhiFA}{\Phi^{\mathrm{FA}}}

\theoremstyle{plain}
\newtheorem{theorem}{Theorem}[section]
\newtheorem{proposition}[theorem]{Proposition}
\newtheorem{lemma}[theorem]{Lemma}
\newtheorem{conjecture}[theorem]{Conjecture}
\theoremstyle{definition}
\newtheorem{definition}[theorem]{Definition}
\newtheorem{construction}[theorem]{Construction}
\theoremstyle{remark}
\newtheorem{remark}[theorem]{Remark}

\title{Equivariance Stratification for Positive Effective Geometry}
\author{}
\date{}

\begin{document}
\maketitle

A positive effective geometry presentation of a BPS positive half becomes a
localisation theorem once the localisation datum has been supplied.
The datum consists of a moduli problem, a group action or reduced
obstruction theory, a proper fixed or quotient locus, and an integration
map. Toric residue formulae, reduced $K3$ classes, CHL twining, and
Borcherds period products occupy different strata. A formula proved in
one stratum requires a fresh datum before propagation to another.

The operadic level is part of the datum. For a CY$_d$ input,
\[
 \Phi_d=\SpCh_{\Sigma_{d-1},C}\circ\PhiFA_d .
\]
At $d\geq 3$ the curve output $A=\Phi_d(\cC)$ is $E_1$-chiral. Braided
$E_2$ structure belongs to centre, double, or module layers such as
$\cZ(\Rep^{E_1}(A))$ after those layers are constructed.

\section{Admissible Equivariance Data}

\begin{definition}[Admissible positive datum]\label{def:admissible-positive-datum}
An admissible positive datum for a CY input $X$ is a quintuple
\[
 \mathfrak L_X=(\Meff(X),G_X,\mathbb E_X,\ell_X,\int_X^{\mathfrak L})
\]
with the following properties.
\begin{enumerate}
\item $\Meff(X)$ is the positive effective moduli stack used for the
Hall or stable-pair construction.
\item $G_X$ is either an algebraic torus, a finite orbifold group, an
arithmetic group on a period domain, or the trivial group together with
a reduced obstruction theory.
\item $\mathbb E_X$ is the perfect obstruction theory, its reduced
cosection quotient, or the automorphic line bundle carrying the
Borcherds product.
\item $\ell_X$ is a proper fixed locus, proper quotient, proper reduced
virtual class, or finite-covolume cusp and divisor datum.
\item $\int_X^{\mathfrak L}$ is the residue, virtual, reduced, orbifold,
or automorphic integration map attached to the preceding data.
\end{enumerate}
Once $\mathfrak L_X$ is fixed,
\[
 \Yp_{\mathfrak L}(X)
 :=
 H^\bullet_{G_X}\bigl(\Meff(X),\mathbb E_X\bigr)_{\ell_X}
\]
is the positive half used in the comparison.
\end{definition}

\begin{proposition}[Four strata]\label{prop:four-equivariance-strata}
The admissible positive data used here fall into four strata.
\begin{enumerate}
\item Toric charts: $G_X=(\C^\times)^d$ or its CY subtorus, and
$\ell_X=\Meff(X)^{G_X}$ is a monomial fixed locus. Integration is
Atiyah-Bott residue.
\item Reduced compact geometry: $G_X$ is absent or finite; the ordinary
virtual class is replaced by a reduced class. For $K3\times E$ the
Oberdieck-Pandharipande theory integrates on the quotient by the elliptic
translation action where the stable-pair quotient and reduced class are
defined.
\item Orbifold and CHL geometry: a finite group acts on the source or on
the elliptic-genus input. Fixed inertia sectors occur for stacky fixed
points. A free CHL quotient contributes through the $g_N$-twined K3
elliptic genus and quotient stable-pair theory; torus residue on
$K3\times E$ is absent.
\item Period-domain geometry: an arithmetic group acts on a Type IV or
Siegel domain. The fixed data are cusps, reflective divisors, multipliers,
and automorphic line bundles; the output is a Borcherds product.
\end{enumerate}
These strata share notation for positive halves, but their integration
maps are different mathematical objects.
\end{proposition}

\section{Toric Fixed-Point Localisation}

Let $U=\C^d$ with torus $T_d=(\C^\times)^d$ and equivariant parameters
$\eps_1,\ldots,\eps_d$. The CY normalisation is
\[
 \eps_1+\cdots+\eps_d=0 .
\]
On $\Hilb^n(\C^d)$ the $T_d$-fixed points are monomial ideals. For
$d=1$ there is one fixed point in each degree. For $d=2$ the fixed
points are Young diagrams. For $d=3$ they are plane partitions and
\[
 \sum_{n\geq0} q^n\,|\Hilb^n(\C^3)^{T_3}|
 =
 \prod_{m\geq1}(1-q^m)^{-m}.
\]

\begin{theorem}[Toric positive half]\label{thm:toric-positive-half}
Let $X$ be a smooth toric CY chart or a toric CY variety whose positive
effective moduli admits a toric fixed-locus decomposition by affine
charts $U_\sigma\simeq \C^d$. Then
\[
 \Yp_{\mathrm{tor}}(X)
 =
 \bigoplus_{\sigma}
 H^\bullet_{T_d}\bigl(\Meff(U_\sigma),\mathbb E_\sigma\bigr)
 \otimes_{\operatorname{Frac}H^\bullet_{T_d}(\mathrm{pt})}
 \operatorname{Frac}H^\bullet_{T_d}(\mathrm{pt}),
\]
and every invariant extracted from this expression is a residue along
monomial fixed ideals. For $X=\C^3$ this gives the positive half
\[
 \CoHA(\C^3)\cong\Yph .
\]
\end{theorem}

\begin{proof}
Atiyah-Bott localisation identifies the equivariant integral with the
sum over connected components of the $T_d$-fixed locus after inverting
the equivariant Euler classes of normal bundles. In the Hilbert-scheme
charts the fixed components are indexed by monomial ideals, hence by the
partitions above. The last statement is the Schiffmann-Vasserot
identification of the cohomological Hall algebra of the Jordan
three-loop potential with the positive affine Yangian half.
\end{proof}

\section{Reduced Compact Geometry}

A compact CY target lacks a toric chart covering its stable-pair theory.
For a K3 surface the holomorphic symplectic form gives a cosection of
the obstruction sheaf, so ordinary virtual counts vanish and the reduced
class is the relevant class. For $S\times E$, with $S$ a K3 surface, the
elliptic curve acts by translation on the $E$ factor. Stable-pair
invariants are reduced and then divided by this translation action on
the locus where the quotient theory is defined.

\begin{proposition}[Reduced $K3\times E$ datum]\label{prop:k3e-reduced-datum}
Let $S$ be a K3 surface and let $X=S\times E$. For a curve class and
stability condition in the Oberdieck-Pandharipande reduced theory, the
positive datum is
\[
 \mathfrak L_X^{\mathrm{red}}
 =
 \bigl(P_n(X,\beta,d)/E,\;1,\;\mathbb E_X^{\mathrm{red}},\;
 [P_n(X,\beta,d)/E]^{\mathrm{red}},\;
 \int_{[P_n(X,\beta,d)/E]^{\mathrm{red}}}\bigr).
\]
The symbol $\C^\times$ in this stratum records the weight of the
reduced obstruction line or holomorphic symplectic form. It records a
line weight, separate from torus actions producing isolated fixed points
on the compact threefold.
\end{proposition}

\begin{proof}
The K3 holomorphic symplectic form gives the cosection used in the
reduced obstruction theory of Maulik-Pandharipande-Thomas and
Kool-Thomas. On $S\times E$, translation by $E$ preserves the
stable-pair problem; Oberdieck-Pandharipande define the reduced
quotient theory by removing this free direction before integration.
Atiyah-Bott denominators are absent from this quotient theory. Such
denominators require an additional local toric chart or finite
fixed-point model.
\end{proof}

\begin{theorem}[The $K3\times E$ invariant separation]
\label{thm:k3e-invariant-separation}
For $X=K3\times E$ the compact and automorphic constants are
\[
\begin{aligned}
 \kcat(X)
 &=\chi(\cO_{K3})\chi(\cO_E)=2\cdot 0=0,\\
 \kch(A_X)
 &=\sum_{q=0}^3(-1)^q h^{0,q}(X)
   =1-1+1-1=0,\\
 \kch^{\mathrm{Heis}}(X)
 &=3,\\
 \kBKM(\Delta_5)
 &=c_1(0)/2=10/2=5,\\
 \kfib(X)
 &=\operatorname{rank}\widetilde H(K3,\Z)=24 .
\end{aligned}
\]
The number $3$ belongs to the Heisenberg-Mukai specialisation. The
number $5$ belongs to the Borcherds denominator. The number $24$
belongs to the K3 Mukai lattice. These three constants lie outside the
compact Hodge supertrace of $K3\times E$.
\end{theorem}

\begin{proof}
K\"unneth gives
$h^{0,\bullet}(K3\times E)=(1,1,1,1)$ and
$\chi(\cO_{K3\times E})=\chi(\cO_{K3})\chi(\cO_E)=2\cdot0=0$. The
Heisenberg-Mukai specialisation reads the K3 contribution $2$ together
with the elliptic Heisenberg line $1$. Borcherds' weight theorem gives
$\operatorname{wt}(\Delta_5)=c_1(0)/2$ for the Gritsenko-Nikulin input
$\Phion=y^{-1}+10+y+O(q)$, hence $\kBKM(\Delta_5)=5$. The Mukai
lattice $\widetilde H(K3,\Z)$ has rank $24$.
\end{proof}

\begin{proposition}[Hall half and denominator recognition]
\label{prop:hall-half-denominator-recognition}
The Hall positive half, the doubled Hall algebra, and the BKM
denominator target are three distinct objects.
\begin{enumerate}
\item For the toric chart $\C^3$, the Hall positive half is
$\CoHA(\C^3)\cong\Yph$.
\item For $K3\times E$, the positive Hall half is the source side of a
Hall-Drinfeld double. The denominator $\Delta_5$ is recognised at the
completed double or target BKM layer after the Hall pairing, radical
quotient, PBW basis, orientation character, and product normalisation
have been fixed.
\item The centre or module layer supplies the braided comparison. It is
additional structure on top of the $E_1$ output at $d=3$.
\end{enumerate}
\end{proposition}

\begin{proof}
A positive Hall half is generated by effective classes and carries the
Hall multiplication. A denominator identity is a statement about a
completed signed product over roots of a generalized Kac-Moody
superalgebra. Passing from the first object to the second requires both
halves, a pairing, a quotient by the radical, a completion, and an
orientation character. The toric identity
$\CoHA(\C^3)\cong\Yph$ therefore stops at the positive-half level,
whereas the $K3\times E$ denominator statement concerns the recognised
BKM target attached to $\Delta_5$.
\end{proof}

\section{Orbifold Inertia and CHL Chambers}

Let a finite group $G$ act on a CY target. For a quotient stack with
fixed sectors, Chen-Ruan cohomology is
\[
 H^\bullet_{\CR}([\Meff/G])
 =
 \bigoplus_{[g]}
 H^{\bullet-2\operatorname{age}(g)}
 \bigl(\Meff^g\bigr)^{C_G(g)} .
\]
If the action is free on the target, the coarse quotient has trivial
geometric inertia beyond the identity sector; CHL sectors then enter
through twined elliptic genera, equivariant sheaf theory, or quotient
stable-pair theory.

\begin{proposition}[CHL input]\label{prop:chl-input}
For $N\in\{1,2,3,4,6\}$ let
\[
 X_N=(K3\times E)/(g_N,t_N)
\]
where $g_N$ is a symplectic K3 automorphism of order $N$ and $t_N$ is
translation by an $N$-torsion point of $E$. The action on $K3\times E$
is free. The automorphic input is the $g_N$-twined weak Jacobi form
\[
 \phi_N=\phi^{(g_N)}_{0,1}(\tau,z),\qquad
 \phi_1=y^{-1}+10+y+O(q),
\]
with constant terms
\[
\begin{array}{c|ccccc}
N & 1 & 2 & 3 & 4 & 6\\
\hline
c_N(0) & 10 & 8 & 6 & 4 & 2\\
\kBKM(\Phi_N) & 5 & 4 & 3 & 2 & 1
\end{array}
\]
in the primitive CHL normalisation.
\end{proposition}

\begin{proof}
The freeness is supplied by the elliptic translation factor. The
twining data are read from the symplectic K3 action on the elliptic
genus. Borcherds' automorphic product theorem gives
\[
 \kBKM(\Phi_N)=\operatorname{wt}(\Phi_N)=c_N(0)/2 .
\]
At $N=1$ the product is the Gritsenko-Nikulin denominator $\Delta_5$.
\end{proof}

\begin{proposition}[Nikulin boundary orders]\label{prop:nikulin-boundary-orders}
The order-indexed Nikulin extension at $N\in\{5,7,8\}$ is separate from
the primitive compact $K3\times E$ CHL ladder. Its automorphic constants
are
\[
\begin{array}{c|ccc}
N & 5 & 7 & 8\\
\hline
c_N(0) & 4 & 2 & 2\\
\kBKM(\Phi_N) & 2 & 1 & 1
\end{array}
\]
whenever the corresponding Borcherds input package has been supplied.
These constants lie outside the compact $K3\times E$ CHL quotient
range.
\end{proposition}

\begin{proof}
For symplectic K3 automorphisms of orders $5,7,8$, the twined Jacobi
input gives the zeroth coefficients $(4,2,2)$ in the order-indexed
Nikulin normalisation. Borcherds' weight theorem gives
$\kBKM(\Phi_N)=c_N(0)/2$, hence $(2,1,1)$. A compact
$K3\times E$ CHL quotient requires a matching finite elliptic
translation construction in the primitive CHL range
$N\in\{1,2,3,4,6\}$. The entries $5,7,8$ therefore require their own
source package before any compact Hall or reduced stable-pair
interpretation is asserted.
\end{proof}

\section{Lattice-Polarised Period Domains}

Let $M\hookrightarrow\Lambda_{K3}$ be a primitive lattice embedding of
signature $(1,r-1)$. The period domain for $M$-polarised K3 surfaces is
the Type IV domain
\[
 \Omega_M
 =
 \{[\omega]\in\Pp(M^\perp\otimes\C):
 (\omega,\omega)=0,\;(\omega,\overline\omega)>0\},
\]
and the moduli stack is an arithmetic quotient
$\Gamma_M\backslash \Omega_M$ after removing the usual discriminant
hyperplanes. Localisation in this stratum is arithmetic: cusps,
reflective divisors, Weyl chambers, and automorphic line bundles replace
torus fixed points.

\begin{proposition}[Period-domain positive datum]\label{prop:period-positive-datum}
A Borcherds lift attached to an $M$-polarised K3 datum defines
\[
 \mathfrak L_M^{\mathrm{per}}
 =
 \bigl(\Gamma_M\backslash\Omega_M,\;\Gamma_M,\;\mathcal L^k\otimes\nu,\;
 \mathrm{div}(\Phi),\;\operatorname{Borch}\bigr),
\]
where $\mathcal L$ is the Hodge line bundle, $\nu$ is the multiplier,
and $\operatorname{Borch}$ is the singular-theta or additive-lift
construction producing the automorphic form $\Phi$. This datum produces
a denominator target. Compact Hall correspondences require additional
source data.
\end{proposition}

\section{Borcherds Weights and Normalisations}

\begin{theorem}[Primitive CHL Borcherds weights]\label{thm:primitive-chl-weights}
For the primitive CHL denominator branch
$N\in\{1,2,3,4,6\}$,
\[
\begin{array}{c|ccccc}
N & 1 & 2 & 3 & 4 & 6\\
\hline
c_N(0) & 10 & 8 & 6 & 4 & 2\\
\kBKM(\Phi_N) & 5 & 4 & 3 & 2 & 1
\end{array}
\]
The $N=1$ form is $\Phi_1=\Delta_5$ and
\[
 D_5=64^{-1}\Delta_5 .
\]
The unnormalised Igusa square is
\[
 \Phi_{10}^{\mathrm{un}}=\Delta_5^2 ,
\]
while the Oberdieck-Pandharipande scalar convention uses
\[
 \Phi_{10}^{\mathrm{OP}}=D_5^2=4096^{-1}\Delta_5^2 .
\]
\end{theorem}

\begin{proof}
Borcherds' theorem gives the weight of the automorphic product as half
the constant term of the input Jacobi form. The primitive CHL twined
inputs have constant terms $(10,8,6,4,2)$ for
$N=(1,2,3,4,6)$, hence the displayed $\kBKM$ row. The coefficient
of $q^0y^0$ in $\phi_1$ is $10$, so the weight of the denominator
$\Delta_5$ is $5$. The monic product divides by the leading theta
coefficient $64$, giving $D_5=64^{-1}\Delta_5$. Squaring removes the
sign character and gives the scalar Igusa branch; the OP scalar uses
the monic product.
\end{proof}

\begin{proposition}[Gritsenko-Clery catalogue]\label{prop:gc-catalogue}
The diagonal-divisor Gritsenko-Clery catalogue is the eight-row table
\[
\begin{array}{c|c|c|c}
\text{row} & (t,N) & \operatorname{wt}F_{t,N} & c_{t,N}(0)\\
\hline
1 & (1,1) & 5 & 10\\
2 & (2,1) & 2 & 4\\
3 & (1,2) & 3 & 6\\
4 & (3,1) & 1 & 2\\
5 & (1,3) & 2 & 4\\
6 & (4,1) & \tfrac12 & 1\\
7 & (1,4) & \tfrac32 & 3\\
8 & (2,2) & 1 & 2
\end{array}
\]
with row-wise formula
\[
 \kBKM(F_{t,N})=\operatorname{wt}F_{t,N}=c_{t,N}(0)/2 .
\]
The common row $(1,1)$ coincides with the primitive CHL $N=1$
denominator.
Comparison with the CHL ladder requires an explicit map of levels,
characters, multipliers, and divisor normalisations.
\end{proposition}

\begin{proof}
Clery-Gritsenko classify genus-two Siegel modular forms with simplest
diagonal divisor on the relevant paramodular groups. Borcherds'
weight formula gives $c_{t,N}(0)/2$ row by row. The primitive CHL ladder
and the diagonal-divisor catalogue are different enumerations after the
common $\Delta_5$ row.
\end{proof}

\section{Admissibility Test}

\begin{proposition}[Admissibility of positive-half formulae]
\label{prop:positive-half-admissibility}
Let $X$ be a CY input. A displayed expression
\[
 \Yp(X)=H^\bullet_G(\Meff(X),\mathbb E)
\]
is a theorem once the admissible datum
$\mathfrak L_X=(\Meff,G,\mathbb E,\ell,\int^{\mathfrak L})$ is supplied.
The following substitutions require additional data:
\begin{enumerate}
\item replacing reduced $K3\times E$ integration by toric residues;
\item reading compact $\kch(A_{K3\times E})=0$ as the Heisenberg
specialisation $3$;
\item reading $\kBKM(\Delta_5)=5$ as a compact or fibre invariant;
\item identifying the primitive CHL constants with the eight
Gritsenko-Clery rows away from $(1,1)$;
\item placing $E_2$ structure on a $d=3$ output before passing to a
centre, double, or module layer.
\end{enumerate}
\end{proposition}

\begin{proof}
Each item violates one component of Definition~\ref{def:admissible-positive-datum}.
Toric residues require isolated torus fixed data. Reduced $K3$ and
$K3\times E$ theories require cosection or quotient integration.
Borcherds weights are automorphic line-bundle weights. The GC rows and
the primitive CHL ladder use different levels and multipliers. The
operadic statement for $d\geq3$ lands in $E_1$.
\end{proof}

\paragraph{Sources.}
Atiyah and Bott 1984, Theorem 7.13. Nakajima 1999. Nekrasov 2002.
Kontsevich and Soibelman 2011. Schiffmann and Vasserot 2013, Theorem
1.1. Maulik, Pandharipande, and Thomas 2010, Theorem 1. Kool and
Thomas 2014. Oberdieck and Pandharipande 2016. Bryan and Oberdieck
2018. Bridgeland, King, and Reid 2001, Theorem 1.2. Chen and Ruan
2004. Nikulin 1979 and 1980. Dolgachev 1996. Gritsenko and Nikulin
1997 and 1998. Gritsenko 1999. Borcherds 1998, Theorem 13.3. Clery and
Gritsenko 2013.

\end{document}
